\section{二进制状态与原生参考组织}
本例选择\textbf{二进制 NOR 存储、本地感测与精确数字归约}。INT8 权重由八个独立数字存储位表示，符号和位权在数字域处理。选择 $W[N,K]$，$K=N=128$、$b_S=b_R=1$ Byte，输入 $B_S=128$ Byte，有效 resident 容量 $16384$ Byte；输出用 $16+\lceil\log_2K\rceil=23$ bit 容器。该尺寸使所选原生页和 sector 完整利用，不是所有技术的共同尺寸要求。

\sid{NOR-01} 为完整二进制随机读周期提供量级，\sid{NOR-02} 为 256-Byte 页、4-KiB sector 和内部 program/erase 提供依据。二者通过明确的本地参考组织互补，不合成为同芯片已测 CIM 宏。NOR-01 MirrorBit 的两个存储位置也不作为本例物理晶体管数的换算依据；机器数据计独立二进制存储位和分配容量。

\sid{NOR-04/05} 是另一种电流精调模式：180 nm ESF1 模拟分类器使用差分浮栅权重，分类器有 101780 cell，784--64--10 网络；100-cell 调谐工作以平均约10/35次脉冲达到3\%/0.3\%误差。其读状态、模型更新和交替擦写与普通 binary 页算法有不同终点。本例不要求均匀 cell 导通电流，不将这些精调脉冲加到 binary 页编程之后。数字归约在正确存储和判决的条件下给出精确整数结果。

\subsection{读分片与擦除域分开}
32 个读分片各保存四个输入行，每行128 Byte，故每片512 Byte、两个256-Byte页。片号 $s=i\bmod32$，片内行 $g=\lfloor i/32\rfloor$；八片共享一个4-KiB sector，sector号与其内Byte地址为
\[
e=\lfloor s/8\rfloor,\qquad a=(s\bmod8)512+128g+o.
\]
全体 $i,o$ 一一覆盖四个 sector 的16384个地址，64个有效页均完整承载逻辑权重。该分段布局要求局部 WL/BL 选择、读偏置缓冲、高压隔离和 page 抑制；它是可检验的工程布局，不是对W25Q内部读端口的反向工程结论。

每片安装128个SA，32片共4096 SA，一批读32个输入项×16个输出×8bit；八个输出组复用这些感测节点。每片一次只选一条本地字线，不要求同一位线同时判读32个cell。32片是同一矩阵的分区，不是独立矩阵副本；更新始终只有一个 native page engine 和一个高压更新域。

\clearpage
\section{有限数字保持与完整输入服务}

新增单个4096-bit tile寄存器，保存一批32项×16输出×8bit权重，捕获后跨八个输入位复用，下一批只能在当前批用完后替换。16条数字通道完成32项位乘／归约；16个23-bit部分和寄存器跨行组保持，输入寄存1024bit，完整输出寄存2944bit。另有原生256-Byte页缓冲用于更新；这些资源不构成全矩阵shadow。

\subsection{读周期的原值与采用值}
\sid{NOR-01} PDF p.85 Table57的speed100给出$t_{RC}\ge100$ ns和$t_{ACC}\le100$ ns，条件为512Mb、$V_{IO}=V_{CC}=2.7$--3.6V、$-55$至$105^\circ$C。Table58的speed120采用120 ns，覆盖128Mb至1Gb、$-55$至$125^\circ$C；p.86 Table59的speed130采用130 ns，$V_{IO}=1.65$V至$V_{CC}$。这些是不同产品速度／电压条件的规格，未作为一个芯片的实测概率分位。

本例采用完整本地二进制读 $t_{rd}=100/120/130$ ns，覆盖局部选择、介质读偏置建立、感测和有效输出。该量级从商品随机周期移植到每片128-SA、32片并行的负载条件；感测宽度和隔离资源明确收费，时序成立仍需电路验证。没有按65nm到28nm尺寸比例缩放介质延迟，也不把外部16-bit数据口或SPI连续带宽作为本地吞吐。

数字保持每物理读批额外捕获一$T_D$。四行组×八输出组共32次读和捕获，八输入位仍需256次数字归约，边界捕获／提交两拍。因此
\[
\boxed{\Delta_S=32t_{rd}+32T_D+256T_D+2T_D.}
\]
典型 $t_{rd}=120$ ns、$T_D=5$ ns，分项为 $3840+160+1280+10=5290$ ns，$\rho=24.20$ MB/s。读侧约73\%是物理读，约24\%是数字归约。各批顺序执行，服务间隔等于完整向量时间；没有额外ADC、$T_I$或重复感测收费。

\subsection{资源生命周期和端口约束}
部分和保留到对应输出组完成全部行组，结果寄存保持到向量提交；tile只覆盖当前32×16权重，释放后才替换。页编程或sector擦除时停止求值，不假设高压选择与读选择可同时使用。读源来自正确感测的二进制权重，跨输入位复用不改变NOR数字状态或原计算精度。

\clearpage
\section{完整持续覆写、有限情景与映射接口}

\sid{NOR-02} PDF p.37规定向预擦为FFh的位置写入1--256Byte；页必须对齐，跨页会回绕。数据输入结束才启动self-timed $t_{PP}$，BUSY恢复0为可接受后续操作的终点。p.39的4-KiB erase同样以BUSY结束；p.66给出页编程typ0.4/max3ms、sector擦除typ45/max400ms。这些完整块已经包含内部高压程序、必要verify和恢复，不再叠加裸脉冲或无依据重试。

本地128-bit接口装256Byte需16拍，命令／完成两拍；erase各付两控制拍。页引擎内部实际cell并行度未公开，保留完整native page服务，不由128-bit接口推成128cell写并行。四个sector均专属该矩阵，无额外无关数据需要保护；任意旧值到新值的持续覆写包括全部4次erase和64次program：
\[
\boxed{T_R=64(t_{PP}+18T_D)+4(t_{SE}+2T_D).}
\]
典型为 $25.6+180+0.0058=205.6058$ ms。有效payload为16384Byte，$\tau=0.07969$ MB/s；erase约占87.5\%，因此数字接口变快对更新能力影响很小。无周期refresh或破坏性restore；维护前和局部长期有效能力相同，寿命另列。

% Generated by scripts/check_nor.py --emit.
\begin{table}[htbp]\centering\small
\begin{tabular}{@{}lrrrrr@{}}\toprule
情景 & $\Delta_S$ ($\mu$s) & $\Delta_R$ (ms) & $\rho$ (MB/s) & $\tau$ (MB/s) & $\mathrm{RI}^{*}$\\\midrule
乐观 & 3.780 & 205.602320 & 33.86 & 0.07969 & 424.9\\
典型 & 5.290 & 205.605800 & 24.2 & 0.07969 & 303.6\\
悲观 & 7.060 & 1792.011600 & 18.13 & 0.009143 & 1983\\
\bottomrule\end{tabular}
\caption{同一 binary 本地感测＋数字归约映射的三组条件配对预算。MB/s 为十进制。每行 $B_S=128$ Byte、$B_R=16384$ Byte；显示位数仅供复算。}\label{03_nor_2d:tab:nor-results}\end{table}

乐观与典型采用原文typ写预算，悲观采用max完整program/erase；相应配对读100/120/130ns和共同数字拍2/5/10ns。三者固定同一阵列、SA、保持和page资源。原文未提供比typ更快的完整写预算，故乐观写吞吐仅因控制开销略高；主范围由完整erase/program条件展开，并非刻意扩大绘图圆。

对一次完整装载后求值$U$个向量，$Q_R=16384$ Byte、$Q_S=128U$ Byte，$\mathrm{RI}=U/128$。典型有
\[
\mathrm{RI}^*=\frac{128}{16384}\frac{205605800}{5290}=303.65,
\qquad U^*=\frac{T_R}{\Delta_S}=38866.88=128\mathrm{RI}^*.
\]
等效每页成本 $T_R/64=3212590.625$ ns为持续覆写的\emph{平均}，不是一笔预擦页program的延迟。Table II使用匹配整个矩阵的$T_R$与向量服务间隔；输出规模、输入共享或更新域改变时需重新映射，不能无条件沿用此$U^*$。

\clearpage
\section{组织对照、解释与证据入口}

如果实现限制使每个512-Byte读分片必须独占4-KiB sector，则仍有4096SA、4096-bit保持和64有效页，但物理分配变128KiB，完整重写需32次erase。未使用位置不增加逻辑payload。
% Generated by scripts/check_nor.py --emit.
\begin{table}[htbp]\centering\small
\begin{tabular}{@{}lrrrrr@{}}\toprule
组织 & 分配容量 & 擦除数 & $\Delta_R$ (ms) & $\tau$ (MB/s) & $\mathrm{RI}^{*}$\\\midrule
主布局 & 16 KiB & 4 & 205.605800 & 0.07969 & 303.6\\
独占 sector & 128 KiB & 32 & 1465.606080 & 0.01118 & 2164\\
\bottomrule\end{tabular}
\caption{保持 32 个读分片、4096 SA 和参考读时间，只改变读分片与擦除域的归属。两行 $\rho$ 相同，物理容量不要求相等。}\label{03_nor_2d:tab:nor-contrast}\end{table}


典型同读侧下，更新能力变为0.01118MB/s、$\mathrm{RI}^*=2164.47$；约7.13倍差异来自擦除域利用率。主布局能否让八个独立读片共用sector高压域，是决定此差异的物理判断；不能从读端口数自动推断擦除域数。该组织变化独立列出，不混入固定配置三情景。

固定典型更新和数字拍，只将移植读预算改变为96/144ns，$\rho=28.31/21.13$ MB/s，$\mathrm{RI}^*=355.22/265.15$。此项仅检查读周期工程桥接，主三情景仍保留完整写侧的typ/max证据。该配置的位置主要由物理随机读、跨位复用后的数字工作和持续sector擦除支配；不能用预擦除有限burst删除erase。

\subsection{原文定位与条件保留}
以下页码均为本地 PDF 页序，从 1 开始。资料保存在 \path{literature/03_nor_2d/}，源标识及 SHA-256 与本地 \path{source_manifest.json} 一致。
\begin{enumerate}
\item \sid{NOR-01}，Infineon，\emph{128 Mb / 256 Mb / 512 Mb / 1 Gb GL-S MIRRORBIT Flash, Parallel, 3.0 V, Military}，002-18741 Rev.~*E，2024-09-04。p.~22 给出内部算法及 1$\to$0 编程/擦除恢复为 1 的语义；p.~45 Table~16 给出带条件的完整更新；pp.~85--87 给出随机与页读时序。65 nm MirrorBit 的物理单元组织并未用于本章一 bit 一数字判决的 cell 面积换算。
\item \sid{NOR-02}，Winbond，\emph{W25Q128JV: 3V 128M-BIT Serial Flash Memory with Dual/Quad SPI}，Revision G，2019-04-08。p.~5 给出 256-Byte 页、16 页/4 KiB sector 和耐久/保持特性；p.~14 是 BUSY 定义；pp.~37--39 是 program/erase 起止；pp.~65--66 是接口与内部时序。p.~66 的注 5 仅为带该注号条目的表征说明，不自动扩展为所有 program/erase 数字的实测条件。
\item \sid{NOR-04}，X.~Guo 等，\emph{Fast, Energy-Efficient, Robust, and Reproducible Mixed-Signal Neuromorphic Classifier Based on Embedded NOR Flash Memory Technology}，IEDM，2017。pp.~1--2 的设计、导入精度和测试条件，p.~3 Figs.~1/2/6 的机制和结构，p.~4 Figs.~9/14 的导入统计与动态响应；55 nm 后续结果属于预测，未用作本章实测时间。
\item \sid{NOR-05}，F.~Merrikh Bayat 等，\emph{Model-Based High-Precision Tuning of NOR Flash Memory Cells for Analog Computing Applications}，DRC，2016。p.~1 是算法和 100-cell 实验；p.~2 Fig.~2 为调谐流程、Fig.~3 为脉冲表征、Figs.~5/6 为精度与脉冲数。与 \sid{NOR-04} 同团队/技术线。
\end{enumerate}


\subsection{复算和外部审阅}
输入、证据、原生维度和资源保存在\path{data/inputs.json}；结果顶层\texttt{native\_configuration}和每情景\texttt{mapping\_interface}分别保存组织与完整装载接口。脚本\path{scripts/check_nor.py}使用共享原生block／full-load API，检查全部16384Byte地址、64页／4sector、完整阶段、保持捕获、单位及$U^*$；默认只读，\texttt{--emit}显式更新结果和表格。

外部审阅重点为128-SA分片的判决与负载、32片并行读隔离、八片共sector的高压erase／page抑制，以及重构后完整商品时序的可迁移性。NOR-02每sector至少$10^5$次program/erase和超过20年保持是原产品规格；每次完整覆写消耗活动sector的一次擦除寿命，不将耐久折入瞬时$\tau$。内部检查不替代这些器件／电路判断。
